Modern supply chains operating under the Physical Internet paradigm face a fundamental signal-quality gap: the urgency declared by operators often diverges from the physically computable real urgency. This leads to over-triage (resource waste) and under-triage (hidden rupture risk). In Supply Chain~4.0 contexts marked by volatility and multi-tier dependencies, no existing digital twin architecture fully corrects this cognitive bias while propagating it consistently across a supplier network.

This report summarises a R\&D internship that proposes a cognitive adequacy score $A(t) \in [0, 100]$ to measure and correct the gap between declared operator urgency and a multi-block, DAG-propagated digital twin real urgency $U_r(t)$. The approach combines pairwise elicitation (AHP, Fuzzy Best-Worst Method) for declared urgency, a noisy-OR aggregation of six operational KPI blocks for real urgency, PROMETHEE~II outranking and systematic shock simulation for network-wide criticality, and a calibrated event engine for keeping the underlying KPI state current. The framework is implemented in Python and operated as a persistent, multi-tenant digital twin exercised through a weekly-cadence serious game, with its adequacy signal validated empirically against recorded operational outcomes.
